\documentclass[10pt]{article}

\usepackage{fancyhdr}
\pagestyle{fancy}
\fancyhf{}
\fancyhead[C]{Working Note: What the Expectations in Proposition 4 and Corollary 1 Are Over}
\fancyfoot[C]{\thepage}
\renewcommand{\headrulewidth}{0.4pt}

\usepackage[margin=2cm]{geometry}
\usepackage[T1]{fontenc}
\usepackage[utf8]{inputenc}
\usepackage{amsmath,amssymb,mathtools}
\usepackage{booktabs}
\usepackage{enumitem}
\usepackage{xcolor}

\setlist{nosep,leftmargin=1.3em}
\newcommand{\E}{\mathbb E}

\begin{document}

\noindent\textbf{Handoff note, 2026-11-24.}
Written because the question ``under the learner's model, what does that mean?'' keeps returning for Proposition 4 (information conservation law) and Corollary 1 (dark-room test) of the RLC paper, \texttt{notes\_joery/rlc/sections/split.tex}.
This note fixes the frame.
Nothing here is in the paper unless Section~\ref{sec:status} says so.

\section{The confusion}
The histories are real.
Actions act in the environment, and outcomes and observations come back from it.
The learner's model does not generate data.
What comes from the model is the \emph{expectation}: the learner's forecast of the histories it will see, made before it sees them.
Saying ``histories generated by the learner's model'' is wrong and is what caused the confusion.

\section{Three laws}
Fix a deployed policy and a run that ends at a time $\tau$.
\begin{itemize}
  \item \textbf{Environment law} $P^{\mathrm{env}}$. The law of the real history $X_\tau$. Unknown to the learner.
  \item \textbf{Hypothesis law} $P_{\hat\theta}$. The law of the history that one hypothesis $\hat\theta$ of the learner's class predicts.
  \item \textbf{Forecast law} $\hat P = \sum_{\hat\theta} \hat b_0(\hat\theta)\, P_{\hat\theta}$. The mixture of the hypothesis laws under the initial belief. This is what the learner expects to see.
\end{itemize}
Every $\E$, every entropy and every mutual information in Proposition 4 and Corollary 1 is taken under the forecast law $\hat P$, including the occupancy that defines $I_{\hat\theta}$, $I_{\hat S}$ and $I_B$.
The belief $\hat b_t$ is a function of the real history, so it is a random quantity under any of the three laws.
The initial belief $\hat b_0$ is fixed.

\section{What the two results state}
Under $\hat P$, and with $\hat b_t$ the exact posterior of the learner's class given the history, the conservation law and the dark-room test are exact.
They are statements about the learner's forecast and its own accounting:
\begin{itemize}
  \item what it expects to learn equals the information it expects its outcomes to carry (Proposition 4),
  \item its expected error on the question ``is this fixed point a dark room'' is bounded both ways by the expected entropy left on that question (Corollary 1).
\end{itemize}
Both need the belief to be the exact posterior.
An approximate or amortised belief can lose entropy without having gathered information, and then neither statement holds even as a forecast.
The paper's guard for this is the tempering constant $c_{\min}$.

\section{When the forecast is about the environment}
\textbf{Condition (C).}
The class contains a hypothesis $\theta^\star$ with $\hat b_0(\theta^\star) > 0$ such that, for every policy the agent can deploy, the law of the observed outcome history under the environment equals $P_{\theta^\star}$.

Under (C) the real histories follow one component of the forecast mixture, and
\begin{equation*}
  P_e = \sum_{\hat\theta} \hat b_0(\hat\theta)\, e(\hat\theta) \;\ge\; \hat b_0(\theta^\star)\, e(\theta^\star),
  \qquad\text{so}\qquad
  e(\theta^\star) \le P_e / \hat b_0(\theta^\star),
\end{equation*}
with $e(\hat\theta)$ the probability that the decision is wrong when $\hat\theta$ is true.
The forecast is right on average over the belief, and for the one environment faced it costs a factor of one over its prior mass.
This is the usual Bayesian assumption that the truth lies in the class.

\textbf{Coarse-graining.}
Condition (C) asks $\theta^\star$ to match the \emph{observed outcome histories}, not the environment's state.
A model that is coarser than the world satisfies (C) as long as it predicts those histories correctly.
The results then hold for the environment at the resolution of the model, and the test answers ``dark room or not, at this resolution''.
The fixed $\iota$ of Theorem 1 (the law of $Z$ given the prediction $Y$ is the same under every policy) is necessary for this, but not sufficient, since the model's dynamics must also be right.

\textbf{Correction to an earlier chat statement.}
I said the corollary holds ``provided the model is consistent in the sense of Theorem 1''.
That is too weak.
The right condition is (C), of which the consistency of Theorem 1 is one part.

\section{When it is not}
\begin{itemize}
  \item \textbf{Inconsistent coarse-graining.} The model merges outcomes that the actions can still choose among, the case below requisite variety. No hypothesis matches the real histories, (C) fails, and the results describe only the forecast. The criterion can lower what the model sees as residual noise and reach a dark room that the test cannot see. Named in the paper as future work.
  \item \textbf{Truth outside the class.} Same failure of (C). The belief concentrates on some hypothesis anyway and can be confident and wrong.
  \item \textbf{Approximate belief.} The accounting itself fails, see above.
  \item \textbf{Dissonance.} The amortised $\hat q$ disagrees with the law the model predicts under $\pi^\star$. Agreed frame: this means the iteration has not converged, so it is outside the corollary, which is stated at a candidate fixed point. The sentence that said so was removed from the main text in a later edit and is not currently on the page.
\end{itemize}

\section{Where the paper stands}
\label{sec:status}
State of the Overleaf at commit \texttt{a8f67e6}.
\begin{center}
\begin{tabular}{p{0.3\linewidth}p{0.62\linewidth}}
\toprule
Place & Does it carry the frame? \\
\midrule
Corollary 1 statement & No. The phrase ``under the learner's model'' was cut in compression. \\
Text after Corollary 1 & No. It speaks of the belief deciding and of floors, with no word on what the expectations are over. \\
Proof of Proposition 4 & Partly. It says expectations are under the learner's own model and that the statement is Bayesian. It also says ``each outcome follows the model's prediction'' and ``the approximated belief'', the first is the confusion of Section 1 and the second contradicts the need for an exact posterior. \\
Proof of Corollary 1 & Yes for the transfer, in the paragraph ``Error in the environment'', which states condition (C) in one clause. \\
Conclusion & Partly. It says the floors are statements of the learner's own model and gives the factor. It says ``miscalibrated or too coarse'', where the precise word is a model that fails (C), since a coarse model that predicts the outcome histories is fine. \\
\bottomrule
\end{tabular}
\end{center}

\noindent\textbf{Smallest set of edits that would put the frame on the page} (proposed, not applied):
\begin{enumerate}
  \item One clause in the hand-over before Corollary 1 or in its statement: expectations are over the histories the learner forecasts under its initial belief.
  \item In the proof of Proposition 4, replace ``each outcome follows the model's prediction'' by the forecast law, and ``approximated belief'' by the posterior of the learner's class.
  \item In the conclusion, replace ``miscalibrated or too coarse'' by a model whose class holds no hypothesis that predicts the outcome histories.
\end{enumerate}

\section{Verification}
The identities and inequalities of Corollary 1 under the forecast law were checked by exact enumeration on 150 random finite models, exact to $10^{-14}$.
The transfer bound is a one-line average and was not run numerically.
The coarse-graining statement follows from condition (C) by definition and was not run numerically either.

\end{document}
